% FILE: content/03_model.tex

\section{Онтологическая структура модели}
\label{sec:model}

Данный раздел излагает формальное ядро Ontology of Continua (OC). Он предоставляет читателю
унифицированную нотацию, аксиомы, определения и структурные ограничения, необходимые для
последующих глав о замыкании, синтезе и практическом применении.

OC описывает континуумы в различных научных областях через единый структурный язык.
Этот язык построен из осей, потенциалов, потоков, порогов, границ, циклов и меры
континуумности. Каждый континуум также вложен в окружающее метапространство.
Все определения в этом разделе формально независимы от предметной области на уровне
фактического словаря. В последующих разделах главы замыкания и верификации утверждения
классифицируются по статусу доказательности и эмпирической поддержке; данный раздел
задал общий языковой каркас для таких проверок.

\subsection{Аксоматический фундамент: уровень \texorpdfstring{$K_0$}{K\_0}}

Уровень \(K_0\) — это чисто структурная подоснова.
Он не является физическим пространством: на нём нет времени, энергии, геометрии или динамики.
Его роль — задавать минимальные условия, при которых любой континуум более высокого уровня
логически возможен.

\paragraph{Спецификация \texorpdfstring{\(K_0\)}{K\_0}}
Континуум \(K_0\) задаётся тройкой
\[
    K_0 = (S,\Delta,\mathcal{C}),
\]
где:
\begin{itemize}
    \item \(S\) — множество состояний;
    \item \(\Delta : S \times S \to \mathbb{R}_{\ge 0}\) — структурная функция различий;
    \item \(\mathcal{C}\) — структурное отношение (или семейство отношений), сохраняющее различимость.
\end{itemize}
На этом уровне отсутствуют параметр времени и оператор эволюции.

\paragraph{Аксиома 0.1 (различимость).}
Для любых \(s_1,s_2 \in S\):
\[
    \Delta(s_1,s_2) = 0 \;\Rightarrow\; s_1 = s_2.
\]
Ненулевое структурное различие — минимальное условие различимости.
Континуум не может существовать без по крайней мере двух различимых состояний.

\paragraph{Аксиома 0.2 (нетривиальный порог \texorpdfstring{$\Theta_0$}{Theta\_0}).}
Существует \(\varepsilon > 0\), такое что
\[
    \forall s_1,s_2\in S,\qquad
    s_1 \neq s_2 \Rightarrow \Delta(s_1,s_2) \ge \varepsilon.
\]
Значение \(\Theta_0 = \varepsilon\) играет роль минимального структурного порога: различия меньше \(\varepsilon\) на уровне \(K_0\) не разрешаются.

\paragraph{Аксиома 0.3 (логическая подоснова).}
\(K_0\) не содержит параметра времени и не обладает оператором динамики.
Он не эволюционирует и сам по себе не порождает уровни выше.
Он задаёт лишь логические условия различимости и существования нетривиальных различий.
Все динамические понятия, а также все представления об энергии, геометрии и причинности
относятся к уровням \(K_1\) и выше.

В частности, никакой оператор, определённый на \(K_0\), не может увеличивать размерность:
различия на уровне \(K_0\) всегда выражаются вдоль существующей структурной оси различимости.
Размерность, связанная с \(K_0\), фиксируется вложенным метапространством \(M_0\).
Здесь \(M_0\) обозначает минимальное метапространство, в котором чисто структурная
подоснова может быть оценена для последующего возникновения \(K_0\to K_1\); это не
физический контейнер и не динамическая среда.

\subsection{Построение уровня \texorpdfstring{$K_1$}{K\_1}}

Уровень \(K_1\) — это наиболее простой подлинный континуум: он вводит время,
одномерную ось и базовую геометрическую структуру.

\paragraph{Спецификация \texorpdfstring{\(K_1\)}{K\_1}}
Континуум \(K_1\) задаётся как
\[
    K_1 = \big(X_1,\tau_1,\Omega_1,A_1,P_1(t),J_1(t),\Theta_1,\partial\Omega_1,C_1,k_1(t)\big),
\]
где:
\begin{itemize}
    \item \(A_1\) — единичная ось (одномерная координата);
    \item \((X_1,\tau_1)\) — топологическое пространство, получаемое из структурных данных \(K_0\), обычно интервал \(X_1 = (a,b)\);
    \item \(\Omega_1\) — пространство допустимых конфигураций на \(X_1\) с подходящими условиями регулярности.
          В данном временно-зависимом аналитическом примере
          \[
              \Omega_1 = C^0(I,H^1(X_1,V_1)) \cap C^1(I,L^2(X_1,V_1)),
          \]
          где \(I\) — временной интервал, а \(V_1\) — гильбертово пространство,
          используемое в примере; частный случай с конечной размерностью —
          конечномерное векторное пространство, применяемое когда локальная модель
          имеет конечное число степеней свободы.
    \item \(P_1(t)\) — функционал, аналогичный потенциальной энергии, на \(\Omega_1\);
    \item \(J_1(t)\) — потоки, получаемые как вариации \(P_1(t)\);
    \item \(\Theta_1\) кодирует минимальные условия для классической устойчивости;
    \item \(\partial\Omega_1\) — граница, где насыщаются пороги устойчивости;
    \item \(C_1\) — множество классических циклов (например, периодических орбит);
    \item \(k_1(t)\) измеряет одномерную континуумность согласно общему определению ниже.
\end{itemize}

\paragraph{Переход \texorpdfstring{$\Psi_{0\to 1}$}{Psi\_0\_to\_1}}
Переход \(K_0 \to K_1\) задаётся оператором
\[
    \Psi_{0\to 1} : (S,\Delta,\mathcal{C}) \mapsto (X_1,\tau_1,A_1,\Omega_1,\Theta_1),
\]
который строит скелетный каркас \(K_1\): энергетический функционал,
семейство потоков, граница, циклы и скаляр континуумности задаются далее в
спецификациях уровня \(K_1\), а не самим этим первым переходом.
Конструкция каркаса включает:
\begin{itemize}
    \item выделение семейства структурно совместимых состояний в \(S\);
    \item наделение их порядком и топологией \((X_1,\tau_1)\);
    \item задание первой оси \(A_1\);
    \item построение допустимого конфигурационного пространства \(\Omega_1\);
    \item индукцию классических порогов \(\Theta_1\).
\end{itemize}
Это первый случай размерностного возникновения: появляется непрерывная ось, которую нельзя
представить внутри чисто структурной подосновы \(K_0\).

\subsection{Общее определение континуума}

Для любого уровня \(K\) в иерархии континуум задаётся кортежем
\[
    K = \big(\Omega(K), A(K), P(t), J(t), \Theta(K), \partial\Omega(K), C(K), k(K,t)\big),
\]
где:
\begin{itemize}
    \item \(\Omega(K)\) — непустое множество допустимых состояний;
    \item \(A(K) = \{A_1,\dots,A_n\}\) — конечное множество осей несовместимых различий;
    \item \(P(t)\) — вектор потенциалов на \(\Omega(K)\);
    \item \(J(t)\) — агрегированный вектор потенциалов и потоков; обновление состояния
          задаётся оператором эволюции \(E\);
    \item \(\Theta(K) = \{\Theta_i\}\) — множество пороговых функций;
    \item \(\partial\Omega(K)\) — граница допустимых состояний;
    \item \(C(K)\) — семейство структурно устойчивых циклов;
    \item \(k(K,t)\) — мера континуумности.
\end{itemize}

Каждый континуум предполагается вложенным в окружающее метапространство \(M\):
\[
    \Omega(K)\subseteq\Omega(M), \qquad A(K)\subseteq A(M).
\]
Метапространство предоставляет дополнительные допустимые состояния и оси, способные
обеспечить будущие размерностные расширения \(K\).

\paragraph{Пространство состояний и граница}
Пороги задаются функциями \(f_i : \overline{\Omega(K)} \to \mathbb{R}\) с условием
\[
    f_i(s) \le 0 \quad\text{для всех } s \in \Omega(K),
\]
а граница определяется как
\[
    \partial\Omega(K) = \{ s \in \overline{\Omega(K)} \mid \exists\, i: f_i(s) = 0 \}.
\]
Такое определение не требует метрики и одинаково применимо к геометрическим,
топологическим, энергетическим, информационным и логическим континуумам.

Граница динамична. Ниже показано правило обновления, а не производная
a на объекте многозначной величины:
\[
    \partial\Omega(K,t+dt) =
    R\big(\partial\Omega(K,t),P(t),J(t),\Theta(K)\big),
\]
которое может сжимать, расширять или ветвить \(\Omega(K)\) в ходе рождения, жизни и смерти.

\subsection{Таксономия порогов}

Каждый континуум обладает структурированным множеством порогов \(\Theta(K)\),
классифицируемых по следующим типам:

\begin{itemize}
    \item \textbf{Пороги существования} \(\Theta_{\mathrm{exist}}\): условия, при которых \(\Omega(K)\neq\emptyset\).
    \item \textbf{Пороги устойчивости} \(\Theta_{\mathrm{stab}}\): условия для ограниченной динамики и недивергентных потоков.
    \item \textbf{Критические пороги} \(\Theta_{\mathrm{crit}}\): гиперповерхности, на которых происходят качественные изменения (фазовые переходы, бифуркации).
    \item \textbf{Размерностные пороги} \(\Theta_{\mathrm{dim}}\): условия, при которых могут возникать новые оси и более высокоразмерные континуумы.
    \item \textbf{Пороги коллапса} \(\Theta_{\mathrm{death}}\): структурные пределы, при которых не остаётся допустимых состояний и \(\Omega(K)\) коллапсирует.
\end{itemize}

Полный ландшафт порогов континуума представляет собой набор неравенств:
\[
    f_{i}^{(\mathrm{exist})}(s) \le 0,\quad
    f_{i}^{(\mathrm{stab})}(s) \le 0,\quad
    f_{i}^{(\mathrm{crit})}(s) \le 0,\quad
    f_{i}^{(\mathrm{dim})}(s) \le 0,\quad
    f_{i}^{(\mathrm{death})}(s) \le 0,
\]
а соответствующие компоненты границы задаются равенствами.

\subsection{Потенциалы, потоки и структурное натяжение}

Потенциалы \(P(t)\) кодируют внутреннюю конфигурацию ограничений и движущих
сил внутри континуума.
Они могут соответствовать энергетическим ландшафтам, химическим концентрациям,
мембранным градиентам, репрезентационным или информационным структурам, а также
институциональному и нормативному давлению.

Агрегированный потоковый вектор \(J(t)\) описывает скорость изменения потенциалов:
\[
    \frac{dP}{dt} = J(t),
\]
а \(J(t)\) раскладывается на структурно различимые компоненты:
\begin{itemize}
    \item \textbf{поддерживающие потоки} \(J_{\mathrm{support}}\), которые поддерживают циклы и стабилизируют \(k(K,t)\);
    \item \textbf{критические потоки} \(J_{\mathrm{critical}}\), которые направляют систему к
          порогам \(\Theta_{\mathrm{crit}}\) и \(\Theta_{\mathrm{dim}}\) или через них;
    \item \textbf{разрушающие потоки} \(J_{\mathrm{kill}}\), которые продвигают систему через
          \(\Theta_{\mathrm{death}}\) и уменьшают \(k(K,t)\).
\end{itemize}

Структурное натяжение \(T(K,t)\) — это функционал потенциалов, осей и градиентов
(схематично \(T=T(P,A,\nabla P)\)).
Оно измеряет степень «натяжения» текущей конфигурации относительно ландшафта порогов.
Размерностные переходы происходят при \(T(K,t) > \Theta_{\mathrm{dim}}(K)\);
коллапс наступает, когда разрушающие потоки совместно со структурным натяжением
приводят систему за пределы \(\Theta_{\mathrm{death}}(K)\).

\subsection{Циклы и континуумность}

Циклы \(C(K)\) — замкнутые траектории в \(\Omega(K)\), которые остаются на конечном
расстоянии от границы и поддерживают континуумность.
На структурном уровне циклы удовлетворяют
\[
    L(C) = \int_C d\ell_K < \infty,\qquad
    S(C) = \min_{s \in C} \delta_{\partial\Omega}(s) > 0,
\]
где \(d\ell_K\) — выбранный структурный элемент длины для реализованного континуума,
а \(\delta_{\partial\Omega}(s)\) — соответствующий запас до границы.
В определении границы выше не требуется метрика; эти оценки циклов вводят
структурную шкалу только при измерении количественного запаса или длины.
В метрических реализациях это может быть индуцированная метрика; в неметрических
реализациях — любой структурно допустимый функционал разделения, равный нулю
ровно на границе и положительный во внутренних структурно допустимых состояниях.

Когда континуум эволюционирует, допустимая область трактуется как зависящая от
времени:
\(\Omega(K,t)\) обозначает пространство состояний, доступное \(K\) в момент \(t\),
а \(\Omega(K)\) — базовую относительную к уровню допустимую область при отсутствии указания
времени.
Когда в более старых формулах встречается \(\Omega(K(t))\), это следует понимать как тот же
объект временного среза: \(\Omega(K,t)\).
Модель вводит следующую единую формулу континуумности:
\[
    k(K,t)
    = \chi_{\Omega}(s_K(t);K,t)\;
      S_{\mathrm{axes}}(K,t)\;
      S_{\mathrm{cycles}}(K,t)\;
      S_{\mathrm{flows}}(K,t)\;
      S_{\mathrm{coh}}(K,t),
\]
где:
\begin{itemize}
    \item \(s_K(t)\) — текущее состояние \(K\) в момент \(t\), а
          \(\chi_{\Omega}(s_K(t);K,t)\) — индикатор существования:
          он равен \(1\), если \(s_K(t)\in\Omega(K,t)\), и \(0\) иначе;
    \item \(S_{\mathrm{axes}}(K,t)\) оценивает эффективность насыщения осей, например,
          отношение эффективного ранга рабочих осей к максимальному возможному рангу для данного уровня;
    \item \(S_{\mathrm{cycles}}(K,t)\) измеряет силу и эффективность устойчивых циклов,
          например взвешенную сумму эффективностей циклов, нормированную на их максимальные значения;
    \item \(S_{\mathrm{flows}}(K,t)\) описывает устойчивость потоков. Для агрегированных
          величин поддерживающих и разрушающих потоков
          \(\Phi_{\mathrm{support}}\) и \(\Phi_{\mathrm{kill}}\) даётся локальное правило:
          \[
              S_{\mathrm{flows}}(K,t)=
              \begin{cases}
              \dfrac{\Phi_{\mathrm{support}}(t)}
                    {\Phi_{\mathrm{support}}(t)+\Phi_{\mathrm{kill}}(t)},
                    & \Phi_{\mathrm{support}}(t)+\Phi_{\mathrm{kill}}(t)>0,\\[0.7em]
              1, & \Phi_{\mathrm{support}}(t)=\Phi_{\mathrm{kill}}(t)=0.
              \end{cases}
          \]
          Если поддержка исчезает, а разрушающий поток остаётся положительным, то
          \(S_{\mathrm{flows}}=0\);
    \item \(S_{\mathrm{coh}}(K,t)\in[0,1]\) кодирует глобальную структурную согласованность,
          включая связность релевантных графов, совместимость порогов и потоков и отсутствие
          противоречивых ограничений.
\end{itemize}
Все мультипликативные факторы в данной формуле нормируются к диапазону \([0,1]\) на локальной модели
перед перемножением.
Оператор обновления континуумности \(U\), формализованный далее в Разделе~the operator semantics chapter,
переводит \(k(K,t)\) на следующий допустимый временной срез с использованием этих факторов.

По построению \(0 \le k(K,t) \le 1\).
Положительное значение \(k(K,t)\) служит индикатором живого состояния.
Оно фиксирует одновременно четыре факта: текущее состояние допустимо,
активные оси поддерживаются, факторы цикла и устойчивости потоков положительны,
а фактор когерентности не коллапсировал. Полный статус живого состояния
требует также условий допустимого состояния, порогов и тождества,
указанных в Разделе~the lifecycle and rebirth chapter.
Смерть соответствует \(k(K,t)\to 0\) в сочетании с коллапсом \(\Omega(K)\).

\subsection{Оператор эволюции}

Эволюция континуума на структурном уровне описывается оператором
\[
    E: K(t) \mapsto K(t+dt),
\]
или в развернутом виде:
\[
    E:\ (\Omega,A,P,J,\Theta,\partial\Omega,C,k)
    \longrightarrow (\Omega',A',P',J',\Theta',\partial\Omega',C',k').
\]
Здесь \(A'\) обозначает множество осей после эволюции того же континуума в расширенном кортеже.
Это унарный оператор эволюции для временного среза \(K\).
Активное метапространство относится к окружению. Когда его требуется указать локально,
оно появляется как параметр с индексом, а не как второй аргумент оператора \(E\),
чтобы арность оператора оставалась фиксированной.

Модель не сводит обновление потоков, проверку порогов, структуру циклов,
движение границы, расширение осей и обновление континуумности в одну
дифференциальную систему. Она рассматривает эти компоненты как связанные части
ограниченной структурной эволюционной схемы с требованиями:

\begin{itemize}
    \item \textbf{Согласованность:} если \(s(t) \in \Omega(K(t))\), то либо
          \(s(t+dt) \in \Omega(K(t+dt))\), либо переход связан с пересечением порога;
    \item \textbf{Соблюдение порогов:} потоки подчиняются неравенствам, закодированным в \(\Theta(K)\),
          за исключением случаев явного пересечения критических или порогов смерти;
    \item \textbf{Монотонность размерности:} если \(A'\) содержит новую ось, не представляющуюся
          комбинацией прежних осей, то \(\dim(K(t+dt)) > \dim(K(t))\); размерность не уменьшается.
\end{itemize}

Динамика продолжается, пока \(\Omega(K(t)) \neq \emptyset\).
Как только \(\Omega(K(t^\ast)) = \emptyset\), континуум умер, и оператор \(E\)
больше не может содержательно действовать на нём.

\subsection{Встраивание в метапространства}

Каждый уровень \(K_x\) вложен в метапространство \(M_x\), которое даёт
дополнительные допустимые состояния и оси.
На структурном уровне это фиксируется условиями
\[
    \Omega(K_x(t)) \subseteq \Omega(M_x(t)), \qquad
    A(K_x) \subseteq A(M_x).
\]
Универсальную эволюционную аксиому можно записать схематически как
\[
    K_x(t+dt) = E_{M_x(t)}\big(K_x(t)\big),
\]
дополнительно требуя: если условие временного вложения не выполняется,
\(\Omega(K_x(t))\not\subseteq\Omega(M_x(t))\), то континуум \(K_x\)
прекращает существование. Если временной параметр ясен, сокращённые обозначения
\(K_x\) и \(M_x\) опускают аргумент \(t\).
Размерностной рост \(K_x\) возможен только если метапространство содержит ось
\(A_{\mathrm{new}}\in A(M_x)\setminus A(K_x)\) и структурное натяжение превышает
соответствующий размерностный порог. Ниже это формулируется как ограничение происхождения
оси: новые оси не могут возникать только из самого \(K_x\).

\subsection{Возникновение континуумов}

Возникновение нового континуума \(K_{x+1}\) из \(K_x\) есть фазовый переход,
индуцированный порогами.
Структурно рождение происходит при выполнении условий:

\begin{enumerate}
    \item Появляется новый класс различий, который нельзя представить через существующие оси \(A(K_x)\).
    \item Структурное натяжение, связанное с этими различиями, удовлетворяет
          \[
              T(K_x,t) > \Theta_{\mathrm{dim}}(K_x).
          \]
    \item Метапространство \(M_x\) содержит по крайней мере одну ось, способную принять новые различия,
          то есть \(A_{\mathrm{new}} \in A(M_x) \setminus A(K_x)\).
    \item Существует непустое множество допустимых состояний \(\Omega(K_{x+1})\), совместимое с новой осью и порогами.
\end{enumerate}

Оператор размерностного рождения
\[
    \Psi_{x\to x+1}: K_x \to K_{x+1}
\]
минимален и необратим: любое ненулевое становление новой оси даёт новый континуум \(K_{x+1}\),
и размерность \(K_{x+1}\) не может вернуться к размерности \(K_x\) без разрушения
\(\Omega(K_{x+1})\).
Локальная модель таким образом реализует размерностную монотонность. Отдельно формулируется
результат об отсутствии спонтанного создания размерности в
Разделе~the spontaneous-dimension theorem discussion.

\subsection{Жизнь континуумов}

Континуум \(K\) считается живым на интервале времени, если
\[
    \Omega(K(t)) \neq \emptyset,\qquad
    k(K,t) > 0,\qquad
    C(K(t)) \neq \emptyset,
\]
и если поддерживающие потоки \(J_{\mathrm{support}}\) доминируют над
разрушающими потоками \(J_{\mathrm{kill}}\) после интегрирования по релевантным циклам.

Жизнь соответствует персистирующему существованию:
\begin{itemize}
    \item устойчивых циклов на конечном структурном расстоянии от границы;
    \item достаточных поддерживающих потоков для поддержания этих циклов;
    \item когерентности между потенциалами, осями и порогами;
    \item контролируемого критического поведения, не пересекающего \(\Theta_{\mathrm{death}}\).
\end{itemize}

Эти условия по-разному интерпретируются на каждом уровне \(K_x\), но структурный
шаблон одинаков от протоцелей до институтов.

\subsection{Смерть континуумов}

Континуум \(K\) умирает в момент \(t^\ast\), когда
\[
    \Omega(K(t^\ast)) = \emptyset,
\]
что влечёт \(\chi_{\Omega}(s_K(t^\ast);K,t^\ast)=0\) и, следовательно, \(k(K,t^\ast)=0\).
Эквивалентно:
\begin{itemize}
    \item все устойчивые циклы исчезают, \(C(K(t^\ast)) = \emptyset\);
    \item нет состояния, одновременно удовлетворяющего всем неравенствам порогов;
    \item любая попытка продолжить динамику нарушила бы хотя бы один порог существования \(\Theta_{\mathrm{exist}}\).
\end{itemize}

\paragraph{Необратимость смерти}
После того как \(\Omega(K(t^\ast)) = \emptyset\), не существует структурного оператора в пределах
того же уровня, который может восстановить непустое \(\Omega(K)\).
Любая кажущаяся «реанимация» соответствовала бы рождению нового континуума \(K'\) со своими
порогами и пространством состояний, но не продолжением исходного \(K\).
Следовательно, смерть структурно необратима.

\subsection{Взаимодействие континуумов}

Два континуума \(K_a\) и \(K_b\) взаимодействуют через оператор
\[
    E_{\mathrm{int}} : (K_a,K_b) \mapsto (K_a',K_b'),
\]
который действует на их совместные потенциалы, потоки, пороги и оси.
Структурно выделяются следующие режимы:

\begin{itemize}
    \item \textbf{Конкуренция:} потоки одного континуума мешают поддержанию циклов в другом;
          как правило, \(k_a(t)\) и \(k_b(t)\) не могут неограниченно возрастать одновременно.
    \item \textbf{Паразитизм:} один континуум использует поддерживающие потоки другого,
          повышая собственную \(k\) за счёт другого континуума;
    \item \textbf{Симбиоз:} поддерживающие потоки сопрягаются так, что оба континуума
          увеличивают континуумность и расширяют допустимые области;
    \item \textbf{Слияние:} оси и потенциалы объединяются в новый континуум
          \(K_{\mathrm{fusion}}\) с объединёнными осями и новым пространством состояний
          \(\Omega_{\mathrm{fusion}}\).
\end{itemize}

Взаимодействие может само по себе индуцировать размерностные переходы,
когда смешанные различия создают ранее неиспользуемые оси и повышают натяжение выше
\(\Theta_{\mathrm{dim}}\).

\subsection{Вертикальная иерархия \texorpdfstring{$K_0$--$K_{12}$}{K0--K12}}

В рамках данной формальной схемы иерархия OC \(K_0,\dots,K_{12}\) сводится к следующему:

\begin{itemize}
    \item \(K_0\): чисто структурная подоснова различимых состояний и различий, без времени и без динамики.
    \item \(K_1\): одномерный классический континуум с энергетическими функционалами и базовыми порогами устойчивости.
    \item \(K_2\): физические континуумы, включая поля, фазовые переходы, перколяцию,
          переходы Berezinskii--Kosterlitz--Thouless и массовые структуры в теории поля.
    \item \(K_3\): химические континуумы, включая реакционные сети, отражательно-автокаталитические
          и энергетически обусловленные структуры, концентрации и параметры среды.
    \item \(K_4\): протоцелльные континуумы с мембранами, осмотическими и кривизными порогами,
          внутренними и внешними градиентами и метаболическими подсистемами.
    \item \(K_5\): ранние нейронные и биоэлектрические континуумы с ионными каналами,
          порогами возбудимости, градиентными потоками и протоспайками.
    \item \(K_6\): когнитивные континуумы с репрезентационными осями, связыванием,
          внутренними моделями, порогами памяти и порогами предсказания.
    \item \(K_7\): социальные континуумы с порогами доверия, институциональными циклами,
          коммуникацией и потоками координации.
    \item \(K_8\): цивилизационные континуумы с крупномасштабной инфраструктурой,
          системными порогами и дальнодействующими циклами.
    \item \(K_9\): теоретические континуумы: теории, парадигмы, онтологии и семьи моделей,
          с порогами согласованности и когерентности.
    \item \(K_{10}\): метатеоретические континуумы, определяющие сравнение, оценку и трансформацию между классами теорий.
    \item \(K_{11}\): рекурсивные метатеоретические континуумы, управляющие семействами операторов,
          металогическим обновлением и рекурсией пространства моделей.
    \item \(K_{12}\): континуумы глобальной когерентности, ограничивающие совместимость между
          семьями структур на уровне \(K_{11}\).
\end{itemize}

Каждый уровень наследует общее определение континуума и добавляет новые оси,
потенциалы, пороги и циклы, специфичные для данного домена.
Эта глава не переописывает все доменно-специфические представления в деталях;
она фиксирует общий структурный каркас, которому все такие реализации должны соответствовать.

\subsection{Итоги}

В этом разделе изложен онтологический каркас OC:
аксиомы \(K_0\), построение \(K_1\), общее определение континуума,
таксономия порогов, роли потенциалов, потоков и структурного натяжения,
а также определение циклов и континуумности.
Также здесь установлены оператор эволюции, встраивание в метапространства,
структурные условия рождения, жизни и смерти, оператор взаимодействия и
вертикальная иерархия \(K_0\)–\(K_{12}\).
Все дальнейшие разработки в Core и в расширениях опираются на эту структуру
как на общий фундамент.
